%% ============================================================
%%  Chapitre 3 — Méthodologie générale
%% ============================================================
\chapter{Méthodologie générale}
\label{ch:methodologie}

\section{Vue d'ensemble du système}

\textit{Scientific Watch Agent} s'articule autour d'un pipeline en cinq
étapes principales, illustré à la figure~\ref{fig:pipeline_overview} :

\begin{enumerate}
  \item \textbf{Recherche} (\textit{Search}) : expansion de la requête et
    récupération d'articles depuis OpenAlex ;
  \item \textbf{Filtrage par qualité} (\textit{Quality Filter}) : scoring
    multi-dimensionnel et sélection des articles les plus pertinents ;
  \item \textbf{Summarisation} (\textit{Summarize}) : génération de résumés
    structurés ou narratifs pour chaque article ;
  \item \textbf{Synthèse} (\textit{Synthesize}) : génération d'une vue
    d'ensemble du domaine à partir des résumés ;
  \item \textbf{Analyse des tendances} (\textit{Trend Analysis}) :
    identification des directions émergentes, des lacunes et des
    perspectives de recherche.
\end{enumerate}

\begin{figure}[htbp]
  \centering
  \includegraphics[width=\textwidth]{figures/pipeline_overview.pdf}
  \caption{Vue d'ensemble du pipeline Scientific Watch Agent. Chaque étape
  correspond à un agent spécialisé dans l'architecture LangGraph.}
  \label{fig:pipeline_overview}
\end{figure}

\section{Source de données : OpenAlex}

\subsection{Présentation d'OpenAlex}

OpenAlex \cite{priem2022openalex} est un index open-source de la littérature
académique mondiale, développé par l'organisation non gouvernementale OurResearch.
Il indexe plus de 250 millions de publications, d'auteurs, d'institutions
et de sources, en exploitant des données publiques telles que Crossref,
PubMed, ORCID et Microsoft Academic Graph.

OpenAlex présente plusieurs avantages déterminants pour notre usage.
Premièrement, son accès est entièrement gratuit via une API REST, contrairement
à des alternatives comme Web of Science ou Scopus qui nécessitent des abonnements
institutionnels onéreux. Deuxièmement, il offre des données structurées riches :
informations sur les auteurs (h-index, affiliations), les venues (classement
par quartile), les citations entrantes et sortantes. Troisièmement, il fournit
une API de filtre avancée permettant des requêtes par concept, domaine,
plage d'années et nombre minimum de citations.

\subsection{Interface OpenAlex dans le système}

Notre module \texttt{src/sources/openalex.py} implémente un client Python
qui encapsule les appels à l'API OpenAlex. Il gère notamment le \textit{cursor
pagination} pour récupérer de larges corpus (jusqu'à plusieurs centaines
d'articles par requête), la déduplication par identifiant OpenAlex,
et le mapping vers le schéma \texttt{Article} défini dans \texttt{src/schemas.py}.

La recherche est effectuée par concepts (\texttt{/concepts}),
en filtrant sur la plage d'années et en triant par nombre de citations
décroissant. Pour chaque article récupéré, les informations d'auteur
(h-index, institution) sont enrichies via des appels secondaires à
\texttt{/authors/\{id\}}, avec une gestion de cache LRU pour limiter
le coût des appels réseau.

\section{Dimensions de qualité : le scorer à six features}
\label{sec:scoring_features}

Nous évaluons la qualité d'un article selon six dimensions complémentaires.
Le score final est une combinaison linéaire pondérée :

\begin{equation}
  \text{score}(a) = \sum_{f \in \mathcal{F}} w_f \cdot s_f(a)
  \label{eq:weighted_score}
\end{equation}

où $\mathcal{F} = \{\text{venue}, \text{authors}, \text{impact},
\text{velocity}, \text{recency}, \text{relevance}\}$,
$s_f(a) \in [0, 1]$ est le sous-score normalisé pour la feature $f$,
et $w_f \geq 0$ avec $\sum_f w_f = 1$.

\subsection{Score de venue (\texttt{venue\_score})}

Le score de venue mesure la qualité de la revue ou conférence dans laquelle
l'article a été publié, selon son classement par quartile Scimago.
Nous définissons un mapping discret : Q1 $\to$ 1.0, Q2 $\to$ 0.7,
Q3 $\to$ 0.4, Q4 $\to$ 0.2, non classé $\to$ 0.1.

Les conférences de premier plan sans classement Scimago (par exemple,
NeurIPS, ICML, CVPR) reçoivent un bonus spécifique via une liste
de venues reconnues maintenue dans \texttt{src/config.py}.

\subsection{Score d'auteurs (\texttt{authors\_score})}

Le score d'auteurs reflète l'autorité scientifique des contributeurs
de l'article, mesurée par l'indice H maximal parmi les auteurs :

\begin{equation}
  s_{\text{authors}}(a) = \frac{\log(1 + h_{\max}(a))}{\log(1 + H_{\text{ref}})}
  \label{eq:authors_score}
\end{equation}

où $h_{\max}(a)$ est l'indice H de l'auteur le plus cité parmi les
contributeurs de $a$, et $H_{\text{ref}} = 50$ est une valeur de référence
normalisatrice. L'application du logarithme atténue l'effet des valeurs
extrêmes (chercheurs à très fort indice H) et évite que ce facteur
n'écrase les autres dimensions.

\subsection{Score d'impact (\texttt{impact\_score})}

Le score d'impact mesure l'influence de l'article via son nombre de
citations rapporté à son âge :

\begin{equation}
  s_{\text{impact}}(a) = \frac{\log(1 + c(a)/\max(1, \text{age}(a)))}
    {\log(1 + C_{\text{ref}})}
  \label{eq:impact_score}
\end{equation}

où $c(a)$ est le nombre de citations, $\text{age}(a)$ est l'ancienneté
de l'article en années, et $C_{\text{ref}} = 20$ est la valeur de référence.

\subsection{Score de vélocité (\texttt{velocity\_score})}

La vélocité de citation mesure la dynamique récente de l'article, c'est-à-dire
la proportion de ses citations accumulées au cours des deux dernières années
par rapport à l'ensemble de ses citations :

\begin{equation}
  s_{\text{velocity}}(a) = \min\left(1, \frac{c_{\text{recent}}(a)}{V_{\text{ref}}}\right)
  \label{eq:velocity_score}
\end{equation}

où $c_{\text{recent}}(a)$ est le nombre de citations des deux dernières années
et $V_{\text{ref}} = 20$ est le seuil de saturation. Cette feature est
particulièrement discriminante pour les domaines d'actualité où les articles
récents à fort impact sont préférables aux travaux fondateurs anciens.

\subsection{Score de récence (\texttt{recency\_score})}

Le score de récence favorise les travaux récents, en appliquant une décroissance
exponentielle en fonction de l'âge de l'article :

\begin{equation}
  s_{\text{recency}}(a) = \exp\left(-\lambda \cdot \text{age}(a)\right)
  \label{eq:recency_score}
\end{equation}

avec une demi-vie de trois ans ($\lambda = \ln(2)/3 \approx 0.231$),
signifiant qu'un article vieux de trois ans voit son score de récence
réduit de moitié par rapport à un article de l'année courante.

\subsection{Score de pertinence (\texttt{relevance\_score})}

Le score de pertinence évalue l'adéquation thématique de l'article
avec le sujet de recherche soumis par l'utilisateur.

\paragraph{V1 — approche lexicale.}
La version lexicale (V1) calcule la correspondance entre les n-grammes
du titre et de l'abstract de l'article et les mots-clés du topic de
recherche, avec un bonus pour les bigrammes.

\paragraph{V2 — approche sémantique.}
La version sémantique (V2) encode le titre et l'abstract de l'article
ainsi que le topic dans un espace vectoriel de dimension 384 via
le modèle \textit{all-MiniLM-L6-v2} \cite{wang2020minilm}, et calcule
la similarité cosinus entre les vecteurs obtenus :

\begin{equation}
  s_{\text{relevance}}^{V2}(a) = \frac{e(a) \cdot e(q)}{\|e(a)\| \cdot \|e(q)\|}
  \label{eq:relevance_v2}
\end{equation}

où $e(a)$ est l'embedding de $a$ et $e(q)$ celui du topic $q$.
Le modèle est chargé localement (sans appel API), mis en cache via
un cache LRU, et son empreinte mémoire est de l'ordre de 80 Mo.

\section{Protocole d'évaluation}

\subsection{Construction des oracles}
\label{sec:oracle_construction}

L'évaluation d'un système de classement de publications nécessite un
oracle de pertinence : une annotation de référence indiquant, pour un
ensemble d'articles, leur niveau de pertinence vis-à-vis du sujet.

\paragraph{Principe de l'oracle.}
Pour chaque domaine $d$, nous récupérons un ensemble de 3 à 5 articles
de revue (\textit{survey papers}) via OpenAlex. Les articles cités dans
la bibliographie d'au moins deux de ces surveys sont considérés comme
\textbf{grade-2} (très pertinents, $r_i = 2$), ceux cités par un seul
survey comme \textbf{grade-1} (pertinents, $r_i = 1$). Les autres articles
du corpus récupéré pour le domaine constituent le bruit de fond
(\textbf{grade-0}, $r_i = 0$).

Cette heuristique de construction est fondée sur l'hypothèse que les
articles régulièrement cités par plusieurs surveys indépendants sont
les références incontournables du domaine — c'est-à-dire les articles
qu'un système de veille idéal devrait placer en tête de classement.

\paragraph{Réparation des oracles (\textit{oracle repair}).}
Certains domaines, notamment les plus récents ou les moins couverts par
des surveys, produisent des oracles avec peu ou pas d'articles grade-2.
Un oracle sans grade-2 rend le calcul de NDCG@15 non informatif (score
identique pour tout classement). Nous avons développé l'utilitaire
\texttt{promote\_oracle\_grade2.py} qui permet, dans ce cas, de promouvoir
manuellement des articles grade-1 sélectionnés par un expert en grade-2,
sur la base de critères de qualité explicites (articles de venues
prestigieuses, fortement cités, directement liés au topic).

\paragraph{Corpus d'évaluation.}
Au total, 19 oracles ont été construits, couvrant 16 domaines d'IA et
3 domaines hors-IA (biologie computationnelle, NLP médical, finance
quantitative). La taille des corpus varie de 90 à 350 articles, avec
un nombre de grade-2 allant de 5 à 21 selon le domaine.

\subsection{Métrique principale : NDCG@15}

Le NDCG@15 (défini à l'équation~\ref{eq:ndcg}) est choisi comme
métrique principale pour les raisons suivantes :

\begin{itemize}
  \item Il prend en compte la \textit{position} dans le classement : placer
    un article très pertinent en premier position est préférable à le
    placer en quinzième ;
  \item Il est normalisé dans $[0, 1]$, facilitant la comparaison entre
    domaines de tailles différentes ;
  \item Il supporte plusieurs niveaux de pertinence (grade-0, 1, 2),
    contrairement aux métriques binaires comme la précision ou le rappel ;
  \item Le seuil $k=15$ correspond à la pratique usuelle de consultation
    des premiers résultats d'une recherche académique.
\end{itemize}

\subsection{Protocole de validation : Leave-One-Out}

Pour l'évaluation du méta-learner (Chapitre~\ref{ch:metalearning}),
nous adoptons une validation croisée \textit{Leave-One-Out} (LOO-CV).
À chaque itération $i \in \{1, \ldots, 19\}$, le domaine $d_i$ est
exclu de l'ensemble d'entraînement, le méta-learner est ajusté sur les
18 domaines restants, et sa performance est évaluée sur $d_i$. Cette
procédure garantit qu'aucune information du domaine de test ne fuite
dans l'entraînement, et produit une estimation non biaisée de la
capacité de généralisation.

\section{Implémentation et choix technologiques}

\subsection{Contrat de communication — Pydantic v2}

Toutes les données structurées échangées entre les agents sont définies
comme des modèles Pydantic v2 \cite{pydantic2024} dans \texttt{src/schemas.py}.
Cette approche présente deux avantages : la validation automatique des types
à chaque transition entre agents, et une documentation implicite du
contrat d'interface qui facilite l'évolution indépendante des agents.

Les schémas principaux sont : \texttt{Article}, \texttt{QualityScore},
\texttt{ArticleSummary}, \texttt{NarrativeSummary}, \texttt{Synthesis},
\texttt{TrendAnalysis}, \texttt{CriticFeedback}, et \texttt{AgentLog}.

\subsection{Gestion des coûts API}

Chaque appel LLM est tracé via \texttt{AgentLog} qui enregistre le nombre
de tokens consommés, le coût en USD et la durée d'exécution. Les optimisations
de coût mises en place incluent : l'utilisation de modèles légers pour
les tâches répétitives (Nebius Qwen3-30B à \$0.14/M tokens pour la
summarisation), le prompt caching Anthropic (90\,\% de réduction sur
les system prompts récurrents), et l'utilisation de Groq pour les tâches
latence-sensibles (inférence Llama à moins de 500\,ms).
